% =========================================================================
% ภาคผนวก ง: มาตรฐานความปลอดภัย ข้อกำหนดไฟฟ้าแรงสูง และการประเมินวงจรชีวิต LCA
% Appendix D: Safety Standards, High-Voltage Codes and LCA Methodology
% =========================================================================

\chapter{มาตรฐานความปลอดภัย ข้อกำหนดไฟฟ้าแรงสูง และการประเมินวงจรชีวิต LCA}
\label{app:safety_and_lca}

ภาคผนวกนี้รวบรวมเกณฑ์ความปลอดภัยสากลสำหรับการทำงานกับระบบพลาสมาความดันบรรยากาศ ข้อกำหนดการป้องกันอันตรายจากไฟฟ้าแรงสูง (High-Voltage Safety Protocols) ในสภาพแวดล้อมการเกษตรที่มีความชื้นสูง และระเบียบวิธีการประเมินวัฏจักรชีวิต (Life Cycle Assessment, LCA) เพื่อประเมินความคุ้มค่าทางเศรษฐศาสตร์และการลดการปลดปล่อยก๊าซเรือนกระจกตามมาตรฐาน ISO 14040/14044

\section{เกณฑ์มาตรฐานการสัมผัสก๊าซพิษและรังสีจากพลาสมา (Occupational Exposure Limits)}
\label{app:sec_oel_limits}

\index{มาตรฐานความปลอดภัย}\index{โอโซน}\index{ไนโตรเจนไดออกไซด์}\index{รังสียูวี}
การคายประจุพลาสมาในอากาศก่อกำเนิดก๊าซโอโซน ($\ce{O3}$) ไนโตรเจนไดออกไซด์ ($\ce{NO2}$) และรังสีอัลตราไวโอเลต (UV-A, UV-B, UV-C) ซึ่งต้องควบคุมให้อยู่ภายใต้เกณฑ์มาตรฐานความปลอดภัยทางอาชีวอนามัย ดังแสดงในตารางที่ \ref{tab:app_oel_limits}

\begin{table}[H]
\centering
\caption{เกณฑ์ขีดจำกัดการสัมผัสสารเคมีและรังสีในสถานที่ทำงาน (OSHA / ACGIH / NIOSH)}
\label{tab:app_oel_limits}
\begin{tabularx}{\textwidth}{l l l X}
\toprule
\textbf{สปีชีส์ / อันตราย} & \textbf{เกณฑ์มาตรฐาน} & \textbf{ค่าขีดจำกัด (Limit)} & \textbf{ผลกระทบทางสรีรวิทยาและข้อกำหนด} \\
\midrule
โอโซน ($\ce{O3}$) & OSHA PEL & $0.10\text{ ppm}$ ($0.2\ \si{\milli\gram\per\meter\cubed}$) & ค่าเฉลี่ยตลอดเวลาการทำงาน 8 ชั่วโมง (TWA) \\
& ACGIH TLV & $0.05\text{ ppm}$ (Heavy work) & งานหนัก ระคายเคืองทางเดินหายใจ \\
& NIOSH IDLH & $5.0\text{ ppm}$ & อันตรายร้ายแรงต่อชีวิตในทันที \\
\midrule
ไนโตรเจนไดออกไซด์ ($\ce{NO2}$) & OSHA PEL (Ceiling) & $5.0\text{ ppm}$ ($9\ \si{\milli\gram\per\meter\cubed}$) & ห้ามสัมผัสเกินค่านี้ ณ เวลาใดเวลาหนึ่ง \\
& ACGIH TLV-TWA & $0.20\text{ ppm}$ & ป้องกันโรคปอดอักเสบเรื้อรัง \\
\midrule
ไนตริกออกไซด์ ($\ce{NO}$) & OSHA PEL-TWA & $25.0\text{ ppm}$ ($30\ \si{\milli\gram\per\meter\cubed}$) & ก๊าซไม่มีสี เป็นสารตั้งต้นของ $\ce{NO2}$ \\
\midrule
รังสีอัลตราไวโอเลต (UV) & ACGIH TLV & $3.0\ \si{\milli\joule\per\centi\meter\squared}$ ที่ $254\text{ nm}$ & ป้องกันการอักเสบของกระจกตา (Photokeratitis) \\
& ความยาวคลื่น & UV-C ($100$--$280\text{ nm}$) & ต้องติดตั้งฝาครอบป้องกันแสงสะท้อนโพลีคาร์บอเนต \\
\midrule
ระดับเสียงการคายประจุ & OSHA Action Level & $85\text{ dBA}$ (8-hr TWA) & เสียงสตรีมเมอร์พัลส์ความถี่สูง ต้องสวมที่ครอบหู \\
\bottomrule
\end{tabularx}
\end{table}

\section{ข้อกำหนดวิศวกรรมความปลอดภัยไฟฟ้าแรงสูงในแปลงเกษตร (Field High-Voltage Safety)}
\label{app:sec_hv_safety}

\index{วิศวกรรมไฟฟ้าแรงสูง}\index{การต่อลงดิน}\index{การป้องกันน้ำและฝุ่น}
เนื่องจากระบบพลาสมาเกษตรต้องทำงานในโรงเรือนหรือแปลงปลูกที่มีความชื้นสูง ($\text{RH} > 80\%$) อุปกรณ์ขับพัลส์ไฟฟ้าต้องเป็นไปตามมาตรฐานความปลอดภัยทางไฟฟ้าระดับสากล (IEEE 510, NFPA 70E, IEC 61010-1) ดังข้อกำหนดต่อไปนี้:

\begin{enumerate}[leftmargin=2.0em, itemsep=4pt]
    \item \textbf{ระดับการป้องกันฝุ่นและน้ำของตู้ควบคุม (Enclosure Ingress Protection):} ตู้จ่ายกำลังไฟฟ้าและแผงวงจรสวิตชิ่ง SiC/GaN ทั้งหมดต้องมีมาตรฐานการป้องกันไม่ต่ำกว่า \textbf{IP65} (กันฝุ่นสนิทและกันน้ำฉีดพ่นแรงดันต่ำได้ทุกทิศทาง)
    \item \textbf{ระบบสายดินป้องกัน (Protective Earth Grounding):} ความต้านทานดินของระบบสายดินขั้วอ้างอิงพลาสมาต้องมีค่าน้อยกว่า \textbf{$5\ \si{\ohm}$} โดยใช้แท่งกราวด์ทองแดงบริสุทธิ์ความยาวไม่น้อยกว่า $2.4\text{ เมตร}$ ตอกลงในดินเปียกชื้น
    \item \textbf{ระบบตัดตอนฉุกเฉินสองชั้น (Dual Interlock Safety System):}
    \begin{itemize}
        \item สวิตช์แม่เหล็ก (Magnetic Reed Sensor) ที่ฝาเปิด-ปิดตู้ปฏิกรณ์: ตัดไฟแรงสูงทันทีภายในเวลา $t < 5\ \si{\micro\second}$ เมื่อเปิดฝา
        \item เซนเซอร์ตรวจจับกระแสรั่วไหลข้ามระนาบ (Arc Flash / Ground Fault): ตัดวงจรไฟฟ้ากระแสสลับเข้าตู้ทันทีเมื่อกระแสรั่วไหล $I_{\text{leak}} > 10\ \si{\milli\ampere}$
    \end{itemize}
    \item \textbf{ระยะห่างฉนวนตามแนวผิว (Creepage and Clearance Distances):} สำหรับแรงดันพัลส์ยอด $10\text{ kV}$ ระยะห่างตามแนวผิวฉนวนเทฟลอน (PTFE) หรือเซรามิกอะลูมินาต้องมีระยะฉนวนไม่น้อยกว่า $25\ \si{\milli\meter}$ ในอากาศแห้ง และไม่น้อยกว่า $50\ \si{\milli\meter}$ ในสภาวะแวดล้อมการเกษตร
\end{enumerate}

\section{ระเบียบวิธีการประเมินวัฏจักรชีวิต (Life Cycle Assessment Methodology)}
\label{app:sec_lca_methods}

\index{การประเมินวัฏจักรชีวิต}\index{คาร์บอนฟุตพริ้นท์}\index{กระบวนการฮาเบอร์-บอช}
การประเมินผลกระทบสิ่งแวดล้อมของการผลิตปุ๋ยไนโตรเจนด้วยพลาสมาเปรียบเทียบกับกระบวนการฮาเบอร์-บอช (Haber-Bosch) ใช้กรอบแนวคิดแบบ \textbf{Cradle-to-Farm-Gate} ขอบเขตของระบบครอบคลุมตั้งแต่วัตถุดิบต้นน้ำ พลังงานไฟฟ้าที่ใช้ในเครื่องปฏิกรณ์ ไปจนถึงการส่งมอบสารละลายปุ๋ยสู่แปลงเกษตร ตารางที่ \ref{tab:app_lca_inventory} แสดงข้อมูลบัญชีรายการสิ่งแวดล้อม (Life Cycle Inventory, LCI)

\begin{table}[H]
\centering
\caption{บัญชีรายการสิ่งแวดล้อมและผลกระทบต่อการปล่อยก๊าซเรือนกระจกต่อ 1 kg ของไนโตรเจนบริสุทธิ์ (N)}
\label{tab:app_lca_inventory}
\begin{tabularx}{\textwidth}{l X X X}
\toprule
\textbf{ตัวชี้วัดผลกระทบสิ่งแวดล้อม} & \textbf{กระบวนการฮาเบอร์-บอช ดั้งเดิม} & \textbf{พลาสมาใช้ไฟจากกริด (Grid)} & \textbf{พลาสมาโซลาร์เซลล์ (Solar PV)} \\
\midrule
การใช้พลังงานปฐมภูมิ (MJ/kg-N) & $35$--$45$ (ก๊าซธรรมชาติ $\ce{CH4}$) & $95$--$140$ (ไฟฟ้ากริด) & $95$--$140$ (พลังงานแสงอาทิตย์) \\
ศักยภาพทำให้โลกร้อน (kg $\ce{CO2}\text{-eq/kg-N}$) & $2.20$--$3.10$ & $4.50$--$6.80$ & \textbf{$0.12$--$0.25$ (ลดลง 95\%)} \\
การใช้น้ำหล่อเย็น ($\si{\liter\per\kilogram}\text{-N}$) & $15$--$25$ & $5$--$8$ & \textbf{$0.5$--$1.2$ (ระบบปิด)} \\
การปล่อยมลพิษขนส่ง (Transport) & สูง (ขนส่งจากโรงงานปิโตรเคมี) & ไม่มี (ผลิต ณ จุดใช้งาน On-site) & ไม่มี (ผลิต ณ จุดใช้งาน On-site) \\
การสูญเสียไนโตรเจนจากการชะล้าง & $30$--$50\%$ (ปุ๋ยเม็ดระเหยเป็นแอมโมเนีย) & $5$--$10\%$ (ปุ๋ยน้ำควบคุมแม่นยำ) & $5$--$10\%$ (ปุ๋ยน้ำควบคุมแม่นยำ) \\
\bottomrule
\end{tabularx}
\end{table}

\begin{equation}
\label{eq:app_lcoe_formula}
\text{ต้นทุนการผลิตสารละลายไนโตรเจนปรับระดับ (LCOE-N):}\quad \text{LCOE}_N = \frac{\sum_{t=0}^T \frac{\text{CAPEX}_t + \text{OPEX}_t}{(1 + r)^t}}{\sum_{t=0}^T \frac{M_{N, t}}{(1 + r)^t}}\quad [\text{บาท/kg-N}]
\end{equation}

เมื่อใช้พลังงานไฟฟ้าจากโซลาร์เซลล์ร่วมกับเครื่องกำเนิดพลาสมาประสิทธิภาพสูง อัตราการคืนทุนของระบบผลิตปุ๋ยในวิสาหกิจชุมชนทุเรียนจันทบุรีเฉลี่ยอยู่ที่ \textbf{2.4 ปี} พร้อมสร้างรายได้เสริมจากคาร์บอนเครดิต (T-VERs) เฉลี่ย $1.8\ \text{ตัน}\ \ce{CO2}\text{-eq}$ ต่อไร่ต่อปี

\clearpage
